% id: 182-370
% book: 개념원리 기하 · 22 평면벡터의 성분
% section: 필수·발전 예제
% figure: none
% answer: $\dfrac{\sqrt{17}}{2}$
% answer_source: 답지
% variant_level: 0 (원본 전사)
% 이 파일은 build.mjs 가 items.json 에서 만든다. 고칠 때는 items.json 을 고친다.
\smash{\raisebox{1.5mm}{\dmkichul{개념원리 기하 182쪽 370번 · 확인체크}}}
$\vec{a}=(3{,}\mkern3mu\mathopen{}-p)$, $\vec{b}=(p+q{,}\mkern3mu\mathopen{}2)$, $\vec{c}=(-3q{,}\mkern3mu\mathopen{}4)$에 대하여 $\vec{c}=2\vec{a}+3\vec{b}$일 때, $|\vec{b}|$를 구하시오.
